\documentclass[11pt,a4paper]{article}

\usepackage{amsmath,amssymb}
\usepackage{graphicx}
\usepackage{booktabs}
\usepackage{hyperref}
\usepackage{xcolor}
\usepackage{tcolorbox}
\usepackage[margin=25mm]{geometry}
\usepackage{xeCJK}
\setCJKmainfont{Hiragino Mincho ProN}
\setCJKsansfont{Hiragino Sans}

\title{「この宇宙」しかあり得ない理由\\
\large 算術的ランドスケープ---高校生のためのガイド}
\author{Wright Brothers Research Program}
\date{2026年4月}

\begin{document}
\maketitle

\begin{tcolorbox}[colback=yellow!5,colframe=orange!60!black,title=このガイドについて]
このガイドは高校生や大学1年生に向けて，Wright Brothers プロジェクトの
論文「算術的ランドスケープ: $L$関数による真空選択とダークエネルギーの
情報コスト」の内容をやさしく解説します．

必要な数学は中学・高校レベル（四則演算，分数，少しの対数）で十分です．
難しい概念はたとえ話で導入します．

\textbf{一番伝えたいこと}: 弦理論は宇宙の``設計図''が $10^{500}$ 通り
あると言いますが，私たちの研究は，物理法則を満たす設計図は
\textbf{たった1通り}しかないことを示します．
その鍵を握るのは「最も小さな素数 $p=2$」です．
\end{tcolorbox}

\tableofcontents

\newpage

%% ============================================================
\section{謎: なぜ「この宇宙」なのか?}
%% ============================================================

\subsection{宇宙は「選ばれた」のか?}

あなたが引っ越しをするとします．不動産屋さんが
「物件は $10^{500}$ 件（1の後にゼロが500個!）あります．
でも地図はありません．カタログもありません．
どれがいいかは住んでみないと分かりません」
と言ったら，どう思いますか?

これが\textbf{弦理論のランドスケープ問題}です．

弦理論（超弦理論）は，素粒子物理学の最有力候補の一つです．
しかしこの理論は，宇宙の ``設計図'' が途方もなく多い
--- 約 $10^{500}$ 通り --- ことを予測します．
それぞれの設計図は ``真空状態'' と呼ばれ，
ダークエネルギーの量や物理定数が少しずつ異なります．

\begin{tcolorbox}[colback=red!5,colframe=red!50!black]
\textbf{ランドスケープ問題}: $10^{500}$ 通りの宇宙が ``可能'' なのに，
なぜ私たちの宇宙はこの特定の姿をしているのか?
弦理論にはこの問いに答える方法がない．
\end{tcolorbox}

弦理論の答えは ``人間原理'' です: 「たまたま我々が存在できる宇宙に
生まれただけだ」というものです．これは科学的な説明というよりは，
あきらめに近い．

\subsection{私たちのアプローチ: たった1軒の家}

Wright Brothers プロジェクトのアプローチは全く異なります．

家のたとえに戻りましょう．私たちの研究は，こう主張します:

\begin{tcolorbox}[colback=green!5,colframe=green!50!black]
\textbf{算術的ランドスケープ}: 物件は $10^{500}$ 件ではなく，
物理法則を満たす家は\textbf{たった1軒}．
引っ越し先を悩む必要はない --- 住める家は1つしかないのだから．
\end{tcolorbox}

では，どうやってそれを示すのか? その鍵が ``$L$関数'' です．

%% ============================================================
\section{$L$関数とは何か?}
%% ============================================================

\subsection{ゼータ関数のおさらい}

以前のガイドで紹介した\textbf{リーマン・ゼータ関数} $\zeta(s)$ を
思い出しましょう:
\begin{equation*}
\zeta(s) = \frac{1}{1^s} + \frac{1}{2^s} + \frac{1}{3^s}
+ \frac{1}{4^s} + \cdots
\end{equation*}

全ての正の整数の $s$ 乗の逆数を足し合わせたものです．

このゼータ関数には ``Euler積'' という別の書き方があります:
\begin{equation*}
\zeta(s) = \frac{1}{1-2^{-s}} \times \frac{1}{1-3^{-s}}
\times \frac{1}{1-5^{-s}} \times \frac{1}{1-7^{-s}} \times \cdots
\end{equation*}

素数 $2, 3, 5, 7, 11, \ldots$ ごとに一つの ``部品'' (Euler因子) がある．
ゼータ関数は，全ての素数の ``部品'' を掛け合わせて組み立てたものです．

前回の結果を一言で: ダークエネルギーを正しく計算するには，
$p=2$ の部品を外す必要がある．$p=2$ を外したゼータ関数を
$\zeta_{\neg 2}(s) = (1 - 2^{-s})\zeta(s)$ と書きます．

\subsection{Dirichlet指標: ``フィルター''}

ここから新しい道具を導入します．

Instagramのフィルターを想像してください．
写真にフィルターをかけると，特定の色が強調されたり，
特定の色が消えたりします．

\textbf{Dirichlet指標}（ディリクレしひょう）$\chi$ は，
数学版の ``フィルター'' です．整数 $1, 2, 3, 4, \ldots$ に対して
``通す'' (値 $1$) か ``ブロックする'' (値 $0$) か
``色を変える'' (値 $-1$ や $i$ など) かを決めるルールです．

たとえば，$p=2$ 用のフィルター（法 $2$ の主指標 $\chi_0$）は:
\begin{itemize}
\item 奇数 ($1, 3, 5, 7, \ldots$) $\to$ 通す (値 $1$)
\item 偶数 ($2, 4, 6, 8, \ldots$) $\to$ ブロックする (値 $0$)
\end{itemize}

偶数を全てブロックするフィルターです．

\subsection{フィルターごとに別の``宇宙''ができる}

各フィルター $\chi$ をゼータ関数にかけると，新しい関数ができます:
\begin{equation*}
L(s, \chi) = \frac{\chi(1)}{1^s} + \frac{\chi(2)}{2^s}
+ \frac{\chi(3)}{3^s} + \frac{\chi(4)}{4^s} + \cdots
\end{equation*}

これが\textbf{Dirichlet $L$関数}です．

ポイント: \textbf{フィルターが違えば $L$関数が違い，
$L$関数が違えばダークエネルギーの値が違う}．

つまり，各フィルターが ``別の宇宙の候補'' を定義するのです．

\begin{tcolorbox}[colback=blue!5,colframe=blue!50!black,title=キーポイント]
\textbf{フィルター $\chi$} $\longrightarrow$
\textbf{$L$関数 $L(s,\chi)$} $\longrightarrow$
\textbf{ダークエネルギー $\Omega_\Lambda(\chi)$} $\longrightarrow$
\textbf{``宇宙の候補''}

異なるフィルターは異なるダークエネルギー値を生み，
異なる ``宇宙の候補'' を定義する．
これが ``算術的ランドスケープ'' --- 数論が定義する宇宙の候補一覧 --- です．
\end{tcolorbox}

%% ============================================================
\section{選別: なぜ $p=2$ だけが生き残るか}
%% ============================================================

\subsection{100\%ルール}

ダークエネルギーの割合 $\Omega_\Lambda$ は，宇宙のエネルギー全体に
占めるダークエネルギーの比率です．

当然ですが，\textbf{100\% を超えることはできません}．
宇宙のエネルギーの120\%がダークエネルギーだったら，
残りの物質のぶんが $-20\%$ になってしまう
--- ありえません．
同様に，ダークエネルギーが $0\%$ 以下（負の値）でも困ります．

つまり:
\begin{equation*}
0 < \Omega_\Lambda < 1 \quad (\text{0\% より大きく 100\% より小さい})
\end{equation*}

これが ``物理的に意味がある宇宙'' の条件です．
この条件でフィルターを選別してみましょう．

\subsection{各フィルターの結果}

各フィルター（素数 $p$に対応する主指標）が与える
ダークエネルギーの値を計算すると:

\begin{center}
\begin{tabular}{cccc}
\toprule
素数 $p$ & ダークエネルギー $\Omega_\Lambda$ & パーセント & 判定 \\
\midrule
$p=2$ & $2\pi/9 \approx 0.698$ & 約 70\% & 合格! \\
$p=3$ & $4\pi/9 \approx 1.396$ & 約 140\% & \textbf{不合格} (100\%超え) \\
$p=5$ & $8\pi/9 \approx 2.793$ & 約 280\% & \textbf{不合格} (100\%超え) \\
$p=7$ & --- & 約 420\%以上 & \textbf{不合格} \\
$\vdots$ & $\vdots$ & $\vdots$ & $\vdots$ \\
\bottomrule
\end{tabular}
\end{center}

\subsection{グラスのたとえ}

\begin{tcolorbox}[colback=cyan!5,colframe=cyan!50!black,title=グラスのたとえ]
コップに水を注ぐことを想像してください．
コップの容量が ``宇宙全体のエネルギー100\%'' に相当します．

\begin{itemize}
\item \textbf{$p=2$ のフィルター}: 水を70\%まで注ぐ
$\to$ 大丈夫! 残りの30\%に物質が入る
\item \textbf{$p=3$ のフィルター}: 水を140\%まで注ごうとする
$\to$ あふれる! 物理的にありえない
\item \textbf{$p=5$ のフィルター}: 水を280\%まで注ごうとする
$\to$ とんでもなくあふれる!
\end{itemize}

$p=2$ 以外のフィルターは，全てコップからあふれてしまいます．
\textbf{あふれない唯一のフィルターが $p=2$} です．
\end{tcolorbox}

\subsection{なぜこうなるのか}

数学的な理由はシンプルです．素数 $p$ のフィルターを使うと，
ダークエネルギーの値は $(p-1)$ に比例して大きくなります:

\begin{itemize}
\item $p=2$: $(p-1) = 1$ $\to$ 最小
\item $p=3$: $(p-1) = 2$ $\to$ もう $100\%$ 超え
\item $p=5$: $(p-1) = 4$ $\to$ さらに大きい
\end{itemize}

$p=2$ は\textbf{最小の素数}であり，$(p-1) = 1$ で
「ぎりぎりちょうどよい」量のダークエネルギーを与える
唯一の素数です．

\begin{tcolorbox}[colback=green!5,colframe=green!50!black,title=結論]
\textbf{物理法則を満たす真空は1つだけ}: 素数 $p=2$ のフィルターが
定義する $\zeta_{\neg 2}(s)$．
この真空はダークエネルギー $\Omega_\Lambda = 2\pi/9 \approx 0.698$
（約70\%）を与える．
\end{tcolorbox}

%% ============================================================
\section{算術的ランドスケープ vs 弦理論ランドスケープ}
%% ============================================================

\subsection{比較表}

弦理論のランドスケープと算術的ランドスケープを比べてみましょう:

\begin{center}
\begin{tabular}{lll}
\toprule
& \textbf{弦理論ランドスケープ} & \textbf{算術的ランドスケープ} \\
\midrule
宇宙の候補数 & $\sim 10^{500}$ 個 & 最大6個 (実質1個) \\
分類の道具 & Calabi--Yau多様体 & Dirichlet $L$関数 \\
& (6次元の複雑な図形) & (素数のフィルター) \\
選択の原理 & 人間原理 & 物理的制約 \\
& (「たまたま」) & ($0<\Omega_\Lambda<1$) \\
$\Omega_\Lambda$ の予測 & できない & $2\pi/9 = 0.698$ \\
実験で反証可能? & 困難 & 可能 (Euclid衛星) \\
\bottomrule
\end{tabular}
\end{center}

\subsection{$10^{500}$ 軒 vs 1軒}

引っ越しのたとえに戻りましょう:

\begin{tcolorbox}[colback=yellow!5,colframe=yellow!50!black,title=家のたとえ]
\textbf{弦理論}: 「物件は $10^{500}$ 軒あります．
地図もカタログもありません．
あなたがたまたまこの家に住んでいるのは，
この家が住みやすかったからでしょう」

\textbf{算術的ランドスケープ}: 「物件は1軒です．
物理法則という ``建築基準法'' を満たす家は1つしかありません．
しかもその家の家賃（ダークエネルギー $\approx 70\%$）まで
正確に分かっています」
\end{tcolorbox}

\subsection{選択肢が少ないほうが良い科学}

「選択肢が多いほうがいいのでは?」と思うかもしれません．
しかし科学では逆です．

\textbf{良い理論} = 予測が具体的で，間違っていることを確かめられる理論．

\begin{itemize}
\item 「明日は雨か晴れか曇りか雪です」$\to$ ほぼ必ず当たるが，何も予測していない
\item 「明日は午後2時に雨が降り始めます」$\to$ 具体的で，外れたらすぐ分かる
\end{itemize}

$10^{500}$ 通りの宇宙を許す理論は，何が観測されても
「その宇宙もランドスケープの中にあります」と言えてしまう．
これは科学的に弱い．

1通りの宇宙しか許さない理論は，観測と1\%でもずれたら終わり．
これは科学的に\textbf{強い}（反証可能性が高い）．

%% ============================================================
\section{半ビット---ダークエネルギーの情報コスト}
%% ============================================================

\subsection{``ビット''って何?}

コンピュータの基本単位は\textbf{ビット} (bit) です．
1ビットは ``2択のうちどちらか'' を知ることに対応します．

\begin{itemize}
\item コインを投げる: 表か裏 $\to$ 結果を知るのに \textbf{1ビット}
\item サイコロを振る: 1から6 $\to$ 結果を知るのに約 \textbf{2.58ビット}
\item じゃんけん: グーチョキパー $\to$ 結果を知るのに約 \textbf{1.58ビット}
\end{itemize}

物理学では，自然対数の底 $e$ を使った ``ナット (nat)'' という単位もよく使います．
$1$ ビット $= \log 2 \approx 0.693$ ナットです．

\subsection{``半ビット'' とは}

では，\textbf{半ビット} $= \frac{1}{2}\log 2 \approx 0.347$ ナット
とは何でしょう?

\begin{tcolorbox}[colback=purple!5,colframe=purple!50!black,title=半ビットのたとえ]
公平なコインを投げると，表と裏はちょうど50:50です．
これは1ビットの情報に相当します．

ではもし，\textbf{コインの片面がほんの少しだけ重い}
ことを知っていたらどうでしょう?
まだ表か裏かは分かりませんが，「完全に五分五分ではない」
という情報は持っている．

この「対称性がわずかに破れている」という知識が，
およそ\textbf{半ビット}の情報に相当します．
\end{tcolorbox}

\subsection{DN境界は半ビットの情報を持つ}

以前のガイドで，宇宙のビッグバン側にはDirichlet境界（D），
未来側にはNeumann境界（N）があることを説明しました．

この ``DN境界'' と，両方がDirichletの ``DD境界'' を比べると，
\textbf{エンタングルメントエントロピー}
（量子的なもつれの度合いを測る量）の差がちょうど:
\begin{equation*}
\Delta S = \frac{1}{2}\log 2 \approx 0.347 \text{ ナット}
= \textbf{半ビット}
\end{equation*}

\begin{tcolorbox}[colback=blue!5,colframe=blue!50!black,title=重要な結果]
\textbf{ダークエネルギーの情報コスト = 半ビット}

ダークエネルギーが存在する理由は，宇宙が ``片方はD，もう片方はN''
という非対称な境界を持つことの ``情報的な代償'' である．
\end{tcolorbox}

\subsection{ダークエネルギー = 時間の非対称性の「代金」}

もう少し直感的に説明しましょう．

\begin{tcolorbox}[colback=orange!5,colframe=orange!50!black,title=ドアのたとえ]
片開きのドアを想像してください．
このドアは一方向にしか開きません（過去 $\to$ 未来）．

両開きのドア（DD境界: 両方Dirichlet）と比べると，
片開きのドア（DN境界: 片方Dirichlet，片方Neumann）は
\textbf{対称性が低い}．

この ``対称性を失った分'' が情報コストとして現れ，
それがダークエネルギーとして宇宙を加速膨張させている．

言い換えると: \textbf{宇宙が「時間に方向がある」
（過去と未来が区別される）ことの代金が，ダークエネルギーである．}
\end{tcolorbox}

コインのたとえでいえば: 宇宙のコインは ``完全に公平'' ではなく，
未来側がほんの少し ``重い''（Neumann条件）．
この非対称性が半ビットの情報を持ち，
それがダークエネルギーとして物理的に現れる．

%% ============================================================
\section{$\log(2)$ が至る所に現れる}
%% ============================================================

$\log 2 = 0.693\ldots$ という数が，この理論の至る所に登場します．
偶然ではありません．

\subsection{$\log 2$ の5つの顔}

\begin{center}
\begin{tabular}{clc}
\toprule
番号 & 登場する場面 & 値 \\
\midrule
1 & 境界エントロピーの差 & $\frac{1}{2}\log 2$ \\
2 & $K_1$レギュレータ（代数的$K$理論の量） & $\log 2$ \\
3 & $p=2$ のEuler因子 $\log\frac{1}{1-2^{-s}}\big|_{s=0}$ & $\log 2$ \\
4 & 境界自由エネルギー差 & $\frac{T}{2}\log 2$ \\
5 & 1ビットの情報のエネルギー（Bekenstein限界） & $k_BT\log 2$ \\
\bottomrule
\end{tabular}
\end{center}

\subsection{全ては $p=2$ に帰着する}

5つの $\log 2$ は，見かけ上は別々の分野（宇宙論，数論，情報理論，
凝縮系物理，代数的トポロジー）に属しています．

しかし，全ての起源をたどると，同じ場所に行き着きます:

\begin{tcolorbox}[colback=green!5,colframe=green!50!black,title=$\log 2$ の起源]
\textbf{全ての $\log 2$ は，最小の素数 $p=2$ と，
それが生み出す $\mathbb{Z}/2$（2で割った余りの世界）から来ている．}

$\mathbb{Z}/2$ とは ``偶数か奇数か'' ``表か裏か'' ``DかNか''
--- 最も基本的な2択の構造です．
\end{tcolorbox}

これは深い意味を持ちます．宇宙の加速膨張（ダークエネルギー）は，
数学の最も単純な構造 --- ``2つに分ける'' --- から来ている可能性がある．

%% ============================================================
\section{ホログラフィー：$p = 2$ は「奥行き」}
%% ============================================================

\subsection{ホログラムのたとえ}

クレジットカードやお札についている\textbf{ホログラムシール}を
見たことがありますか?
あの薄い2次元のシールを傾けると，3次元の像が浮かび上がります．
2次元の表面に3次元の情報が``畳み込まれている''のです．

実は，私たちの研究で見つかった $p = 2$ の構造は，
数学版のホログラムです．

\subsection{数学の``次元''}

数論には``エタールコホモロジー次元''という，
空間の``奥行き''を測る指標があります．

\begin{center}
\begin{tabular}{lcc}
\toprule
数学的対象 & 意味 & ``次元'' \\
\midrule
$\mathrm{Spec}(\mathbb{Z})$ & 全ての素数を含む空間 & 3 \\
$\mathrm{Spec}(\mathbb{Z}[1/2])$ & $p=2$ を除いた空間 & 2 \\
\bottomrule
\end{tabular}
\end{center}

\begin{tcolorbox}[colback=purple!5,colframe=purple!50!black,title=ポイント]
$p=2$ を取り除くと，数学的な``次元''が \textbf{3から2に下がる}．

つまり，$p=2$ こそが ``3次元目''を担っていた --- ホログラムでいう
``奥行き''の方向だったのです．
\end{tcolorbox}

\subsection{ホログラムの構造}

ホログラムシールのたとえに戻りましょう:

\begin{itemize}
\item \textbf{2次元のシール}（表面）= $\mathrm{Spec}(\mathbb{Z}[1/2])$
（$p=2$ 以外の素数の世界，次元 $= 2$）
\item \textbf{奥行き方向}（3次元目）= $p = 2$ の方向
\item \textbf{3次元の像}（全体）= $\mathrm{Spec}(\mathbb{Z})$
（全ての素数の世界，次元 $= 3$）
\end{itemize}

2次元の ``境界''に，$p = 2$ という``奥行き''が加わることで，
3次元の全体像が復元される --- これが算術ホログラフィーです．

DN真空（$p = 2$ を除去する操作）とは，
ホログラムから\textbf{奥行き方向を切り取る}ことに相当します．
その結果，次元が下がり，失われた奥行きの``代償''が
ダークエネルギーとして現れるのです．

\subsection{物理学者はすでに知っていた}

実は，物理学にはすでに似た話があります．

\textbf{ブラックホール情報パラドックス}:
ブラックホールに落ちた情報は消えてしまうのか?
物理学者の答えは「消えない --- ブラックホールの\emph{表面}（2次元）に
全ての情報が記録されている」というものです．
これが\textbf{AdS/CFT対応}と呼ばれる理論で，
3次元の物理を2次元の表面で完全に記述できることを示しています．

\begin{tcolorbox}[colback=blue!5,colframe=blue!50!black,title=AdS/CFTと算術版の比較]
\textbf{物理学のAdS/CFT}: 3次元の重力 = 2次元の表面の理論

\textbf{算術版}: $\mathrm{Spec}(\mathbb{Z})$（次元3）=
$\mathrm{Spec}(\mathbb{Z}[1/2])$（次元2）+ $p=2$ の奥行き

私たちの発見は，物理学で知られていたホログラフィーの
\textbf{算術版}（数論版）なのです．
\end{tcolorbox}

\subsection{4次元時空の分解}

最終的に，私たちの宇宙の4次元時空は次のように分解されます:

\begin{equation*}
\underbrace{4\text{次元}}_{\text{時空}}
= \underbrace{2\text{次元}}_{\text{表面（境界）}}
+ \underbrace{1\text{次元}}_{p=2 \text{ の奥行き}}
+ \underbrace{1\text{次元}}_{\text{時間}}
\end{equation*}

\begin{tcolorbox}[colback=green!5,colframe=green!50!black,title=この章のまとめ]
\begin{center}
\large
\textbf{宇宙は算術的ホログラムであり，\\
$p = 2$ はそれに奥行きを与える方向である．}
\end{center}

\vspace{0.5em}

ホログラムシールが2次元なのに3次元に見えるように，
$p=2$ 以外の素数が作る2次元の ``表面'' に
$p=2$ の ``奥行き'' が加わることで，
私たちの宇宙の3次元空間（+1次元の時間）が構成される．

DN真空で $p=2$ を取り除くことは，
ホログラムから奥行きを消すことに等しい．
その ``代償'' がダークエネルギーである．
\end{tcolorbox}

%% ============================================================
\section{何を実験で検証できるか?}
%% ============================================================

理論がどんなに美しくても，実験で確かめられなければ科学ではありません．
算術的ランドスケープの予測は，具体的に検証可能です．

\subsection{Euclid衛星 (2028年): $\Omega_\Lambda$ の精密測定}

\begin{tcolorbox}[colback=blue!5,colframe=blue!50!black]
\textbf{予測}: $\Omega_\Lambda = \frac{2\pi}{9} = 0.6981\ldots$

\textbf{現在の観測値}: $\Omega_\Lambda = 0.685 \pm 0.007$
（Planck衛星 2018年）

\textbf{今後}: ESAのEuclid衛星は2028年までに $0.5\%$ の精度で
$\Omega_\Lambda$ を測定する予定です．
$0.5\%$ の精度は $\pm 0.003$ 程度であり，
$0.698$ と $0.685$ を区別できます．
\end{tcolorbox}

もし Euclid の測定値が $0.698$ に近ければ，この理論は大きな支持を得ます．
逆に $0.698$ から大きくずれれば，理論は反証されます．

\subsection{マルチバースのシグネチャ: 見つかったら理論は終わり}

弦理論のランドスケープでは，``バブル宇宙'' 同士が衝突することがあり，
その痕跡がCMB（宇宙マイクロ波背景放射）に残る可能性があります．

算術的ランドスケープでは，物理的真空は1つだけなので
マルチバースは存在しません．

\begin{tcolorbox}[colback=red!5,colframe=red!50!black]
\textbf{反証条件}: もしCMBに ``バブル衝突'' のシグネチャが見つかったら，
算術的ランドスケープは\textbf{反証される}．
マルチバースが存在するなら，``宇宙は1つ'' という主張は間違いだからです．

これは良いことです --- 間違いを確かめられることが科学の強さです．
\end{tcolorbox}

\subsection{実験室での検証: 凝縮系物理}

半ビット $\frac{1}{2}\log 2$ の境界エントロピーは，
宇宙論だけでなく，実験室の物質でも現れます．

\begin{itemize}
\item \textbf{Majoranaフェルミオン}: 超伝導体の端に現れる特殊な粒子．
端のエンタングルメントエントロピーがちょうど $\frac{1}{2}\log 2$
であることが，理論的にも実験的にも確認されています．

\item \textbf{トポロジカル絶縁体}: DN型の境界条件を持つ物質系で，
境界エントロピーが $\frac{1}{2}\log 2$ になることが
実際に観測されています．

\item \textbf{量子ドット系}: 微小な半導体素子で，
異なる境界条件の間のエントロピー差を直接測定できます．
\end{itemize}

つまり，半ビットという量は ``宇宙論の空想'' ではなく，
\textbf{実験室で既に確認されている物理量}です．

%% ============================================================
\section{正直な限界}
%% ============================================================

この研究には，重要な限界があります．
科学では「何が分からないか」を正直に言うことが大切です．

\begin{tcolorbox}[colback=gray!5,colframe=gray!50!black,title=この理論の限界]
\begin{enumerate}
\item \textbf{``なぜ $L$関数なのか'' は完全には説明されていない}

「真空状態は $L$関数で分類される」というのは\emph{提案}であり，
もっと根本的な理論から \emph{証明} されたわけではありません．
動機付けはありますが，確定ではありません．

\item \textbf{正規化係数 $8\pi/3$ の起源}

ダークエネルギーの計算式に出てくる係数 $8\pi/3$ は
ホログラフィック原理から仮定されています．
なぜこの値なのかの完全な導出はまだありません．

\item \textbf{半ビットの計算は簡略化されたモデルに基づく}

境界エントロピー $\frac{1}{2}\log 2$ は
自由場（相互作用のない最も単純な場）で計算されたものです．
実際の宇宙にはもっと多くの場があり，
補正が入る可能性があります．

\item \textbf{非主指標の排除は完全な証明ではない}

$q \leq 100$ までの計算機による確認はありますが，
全ての導手（全てのフィルター）に対する
数学的な証明はまだありません．

\item \textbf{ダークマターは説明できない}

この枠組みはダークエネルギー ($\Omega_\Lambda$) のみを予測し，
ダークマターやバリオン物質の量は予測しません．
この限界は確認された否定的結果です．
\end{enumerate}
\end{tcolorbox}

これらの限界は，理論が ``未完成'' であることを意味しますが，
``間違い'' であることを意味しません．
2028年のEuclid衛星の結果が，最終的な判定を下すでしょう．

%% ============================================================
\section{高校生のみなさんへ---宇宙は最も単純な算術で書かれているかもしれない}
%% ============================================================

\subsection{物語のまとめ}

この研究が伝えようとしていることを，もう一度まとめます:

\begin{enumerate}
\item 弦理論は $10^{500}$ 通りの宇宙を許すが，
どれが ``正解'' かを選べない（ランドスケープ問題）

\item 私たちは ``フィルター''（Dirichlet指標）を使って
宇宙の候補を整理した（算術的ランドスケープ）

\item 「ダークエネルギーが $0\%$ より大きく $100\%$ より小さい」
という当たり前の条件だけで，
候補はたった1つに絞られた（$p=2$ のフィルター）

\item その唯一の宇宙のダークエネルギーは $\Omega_\Lambda = 2\pi/9 \approx 70\%$
--- 実際の観測値 ($\approx 68.5\%$) に非常に近い

\item ダークエネルギーの存在理由は ``半ビットの情報コスト''
--- 時間に方向があることの代金

\item 至る所に現れる $\log 2$ は，全て最小の素数 $p=2$ から来ている
\end{enumerate}

\subsection{最も単純な算術}

$p=2$ は最小の素数です．$\mathbb{Z}/2$（偶数か奇数か）は
最も単純な数学的構造の一つです．

もしこの理論が正しければ，宇宙の加速膨張という壮大な現象は，
\textbf{最も小さな素数が生み出す，最も単純な算術的構造} から
生じていることになります．

\begin{tcolorbox}[colback=green!5,colframe=green!50!black,title=メッセージ]
物理学の歴史は，自然が驚くほど単純な数学で記述できることを
繰り返し示してきました．

ニュートンの重力法則は $F = GMm/r^2$ という短い式．
アインシュタインの $E = mc^2$ は中学生でも書ける．
マクスウェルの電磁気学は4つの方程式で全てを説明する．

算術的ランドスケープは，この伝統の延長線上にあります:

\begin{center}
\large
\textbf{宇宙は，最も単純な算術で書かれているかもしれない．}
\end{center}

その ``最も単純な算術'' とは，$2$ で割った余りの世界 --- $\mathbb{Z}/2$ ---
であり，それが生み出すのは:
\begin{itemize}
\item 宇宙のダークエネルギー ($\Omega_\Lambda = 2\pi/9$)
\item 半ビットの情報コスト ($\Delta S = \frac{1}{2}\log 2$)
\item 時間の矢の代金
\end{itemize}

2028年，Euclid衛星がこの物語の結末を教えてくれるでしょう．
\end{tcolorbox}

\subsection{あなたにできること}

\begin{itemize}
\item \textbf{数学を学ぶ}: 高校の数学（特に数列と対数）は，
ゼータ関数と $L$関数を理解する第一歩です

\item \textbf{批判的に考える}: この研究にも限界があることを忘れないでください．
「本当にそうなの?」と疑問を持つことが科学の出発点です

\item \textbf{2028年を待つ}: Euclid衛星の結果は公開されます．
自分の目で確かめることができます

\item \textbf{探求する}: Wright Brothers プロジェクトの全ての論文は
公開されています．興味があれば，もっと深く読んでみてください
\end{itemize}

\vspace{1em}

\begin{center}
\textit{科学の最前線は，実はあなたのすぐそばにある．}
\end{center}

\end{document}
